\section{The reporting deficit and what would close it}
\label{sec:open}

Two things follow from the measurements above: a change the field can make now at no research cost, and
three questions that need experiments nobody has run.

\subsection{A reporting checklist, taken from where the silence is}
\label{subsec:checklist}
\label{subsec:record}

The silence is concentrated enough to be addressed by a list. The twelve dimensions with the highest
weighted \texttt{not\_reported} rate all sit above 76\%, and nine of the twelve describe how a harness is
bounded or recovered; tool count, protocol standardization, and context compaction are the exceptions.
Each becomes a minimum statement: Table~\ref{tab:checklist} gives the five most silent, and
Supplement~S6 all twelve. Each item is
one sentence or configuration line, needs no experiment, and is checkable against the released
artifact; that separates a reporting item from a recommendation.

A negative answer satisfies an item as fully as a positive one, which makes the list cheap to adopt.
``The agent runs in the host shell with no isolation and no network restriction'' answers two rows, and
our scheme codes it as a value rather than as silence. Without such a sentence, neither a reader nor our
coder can tell an unbounded system from an undocumented one. A general call to state the harness exists
\citep{zhang2026harnessreporting}. The audit adds which items go unstated and how often, weighted to the
census, and the list is short because the measurement made it short.

The same gap sits one level up: no board in our harvest of 4,887 leaderboard entries records which
harness commit produced a score. Answering the twelve items would not make harnesses comparable. It
would let a reader see how a system is bounded, and give a design analysis dimensions stated often
enough to use.

\begin{table}[t]
\caption{The five items of the reporting checklist with the highest weighted silence rate; all twelve
are in Supplement~S6. \emph{Silent} is the weighted census share of systems whose sources say nothing
about the dimension (design standard errors in Table~\ref{tab:landscape_layers}), an upper bound
(\S\ref{subsec:humanaudit}). The rates measure
documents, not systems.}
\label{tab:checklist}
\footnotesize
\begin{tabular}{llrp{0.50\linewidth}}
\toprule
Layer & Dimension & Silent & Minimum a paper or repository should state \\
\midrule
G & \texttt{network\_policy} & 93.3\% & Whether the agent can reach the network during a run, and if so what is reachable: open egress, an allowlist, or nothing. \\
H & \texttt{replayability} & 90.6\% & Whether a finished run can be replayed from stored artifacts, and whether fully, from inputs only, or not at all. \\
G & \texttt{filesystem\_access} & 87.5\% & The writable scope: the whole host, one working directory, a copy-on-write overlay, or read-only. \\
E & \texttt{rollback} & 86.4\% & What happens to partial edits when a step fails: snapshot restore, version-control revert, or nothing. \\
D & \texttt{state\_persistence} & 86.2\% & Whether a run can be resumed after interruption, and from what: full state, a checkpoint, or nothing. \\
\bottomrule
\end{tabular}
\end{table}

\subsection{Three open questions}
\label{subsec:openq}

\textbf{Does self-verification help at equal compute?} The pooled ablation effect is an upper bound,
\rqSVPooled{} relative on the completed corpus-wide harvest before its bias discount and about
\rqSVDiscounted{} after it, and its prediction interval, \rqSVPI, includes zero. It is also no longer
distinguished: most poolable dimensions show a gain of similar size, so published ablations cannot say
whether verification matters more than any other component. Verification issues extra model calls,
and none of the 42 coded-set papers that ablate a component runs a compute-matched control, so the
ablations cannot separate design from compute. A three-arm
study would: one attempt; verify-and-repair up to $k$ calls; and the same calls, per instance, spent on
unguided retries. We froze that design on 2026-09-24 and filed it, with a dated revision, as protocol
amendments, now accepted at OSF; only pilot data exist, and no available suite supplies the 93
in-band instances it needs.

\textbf{Do the components that bound an agent matter?} Almost nobody has published a number. On the
coded-set harvest, none of the eleven dimensions of layers F, G, and H carried a usable ablation contrast
across the 42 papers that ablate anything. The corpus-wide harvest reaches four of them, where
termination condition pools to a precise null and guardrails to an interval spanning zero; of the
\rqPreRuleGContrasts{} contrasts it first attributed to layer G, \rqPreRuleGRemapped{} removed
code-execution feedback or file-system tools, not a sandbox or a permission policy, and the other
\rqPreRuleGDropped{} were unreadable or named no sandbox mechanism. Under the
fixed mapping rules of Section~\ref{sec:outcomes}, \rqZeroDimCount{} of the
\rqAblatableDimCount{} ablatable dimensions carry none (\rqZeroDimList). That is an absence of evidence about the literature, not evidence that the components are
harmless to remove.
The nearest controlled measurement varies scaffold architecture rather than a bound and gives the
model no tools: over 62,808 scored evaluations, scaffold architecture explains 0.4\% of the variation
in measured safety and benchmark choice 19.3\% \citep{safetyscaffolding2026}.
The closing study is small: one harness, one task suite, and a bound varied one level
at a time. Its tasks should include ones the bound exists to block, so the bound's cost and the
failure it prevents appear together.

\textbf{Do design families exist in a better-reported corpus?} Ours does not contain them: the best
eligible clustering reached a mean silhouette of 0.120 against a pre-stated floor of 0.25. That result
covers only the half of the corpus reporting enough to be clustered, and \S\ref{subsec:construct} gives a
second explanation our cells cannot rule out. The identical procedure, run on a corpus that answers
the twelve-item checklist as a condition of inclusion and coded with a loop-driver axis, would separate
the two.
